% 中文示例：需要 ctex 宏包，并用 xelatex（或 lualatex）编译。
%
% BasicTeX 用户请先安装：
%   sudo tlmgr update --self
%   sudo tlmgr install ctex xeCJK enumitem
%
% 测试：
%   ./src/convert.sh both examples/demo-zh.tex
%   ./src/preview.sh    examples/demo-zh.tex
\documentclass[11pt]{article}
\usepackage{amsmath,amssymb}
\usepackage{geometry}
\geometry{margin=2.5cm}
\usepackage{ctex}          % 中文支持（务必用 xelatex 编译）

\title{latex-convert · 中文演示}
\author{boff868}
\date{\today}

\begin{document}
\maketitle

\section{这是什么}
这是 \texttt{latex-convert} 的中文示例文稿。把它拖到 App 图标上，
会自动在源文件同目录生成同名的 PDF 和 Word 文件；
选择「预览」则会在浏览器里打开可切换「排版预览 / 源码」的阅读页。

\section{公式与中文}
行内公式 $E = mc^2$，行间公式：
\begin{equation}
  \int_{-\infty}^{\infty} e^{-x^2}\,\mathrm{d}x = \sqrt{\pi}.
\end{equation}

中文排版需要 \texttt{ctex} 宏包，并使用 \texttt{xelatex} 引擎编译；
本工具的引擎查找顺序为
\texttt{latexmk → xelatex → pdflatex → lualatex}。

\section{列表}
\begin{itemize}
  \item Word：由 \texttt{pandoc} 转换；
  \item PDF：由本地 LaTeX 引擎编译（默认两遍，保证交叉引用正确）；
  \item 预览：由 \texttt{preview.sh} 生成 HTML 阅读页。
\end{itemize}

\end{document}
